% ============================================================
%  PRINCIPIOS WCAG 2.2 Y EVIDENCIAS --- VitaPrevent
%  Análisis POUR con evidencias de código y pruebas
%  Compilar con: pdflatex principios_wcag_evidencias.tex (2 veces)
% ============================================================
\documentclass[12pt, a4paper]{article}

\usepackage[utf8]{inputenc}
\usepackage[T1]{fontenc}
\usepackage[spanish, es-noshorthands]{babel}
\usepackage[top=2.5cm, bottom=2.5cm, left=3cm, right=2.5cm]{geometry}
\usepackage[dvipsnames, table]{xcolor}

\definecolor{navydeep}{HTML}{0D1B2A}
\definecolor{royalblue}{HTML}{1565C0}
\definecolor{lightblue}{HTML}{90CAF9}
\definecolor{successgreen}{HTML}{16A34A}
\definecolor{warningyellow}{HTML}{D97706}
\definecolor{errored}{HTML}{DC2626}
\definecolor{graymid}{HTML}{6B7280}
\definecolor{purplemod}{HTML}{7C3AED}

\usepackage{booktabs}
\usepackage{longtable}
\usepackage{array}
\usepackage{multirow}
\usepackage{makecell}
\renewcommand{\arraystretch}{1.4}
\usepackage{enumitem}
\setlist[itemize]{leftmargin=1.4em, itemsep=3pt}
\usepackage{listings}
\lstdefinestyle{wcagcode}{
  backgroundcolor=\color{royalblue!6},
  basicstyle=\ttfamily\footnotesize,
  keywordstyle=\color{royalblue}\bfseries,
  stringstyle=\color{successgreen},
  commentstyle=\color{graymid}\itshape,
  numberstyle=\tiny\color{graymid},
  numbers=left, numbersep=8pt,
  breaklines=true, frame=single,
  rulecolor=\color{lightblue!50},
  xleftmargin=12pt, xrightmargin=4pt,
  tabsize=2, showstringspaces=false,
  extendedchars=true,
  literate=
    {á}{{\'a}}1 {é}{{\'e}}1 {í}{{\'i}}1 {ó}{{\'o}}1 {ú}{{\'u}}1
    {Á}{{\'A}}1 {É}{{\'E}}1 {Í}{{\'I}}1 {Ó}{{\'O}}1 {Ú}{{\'U}}1
    {ñ}{{\~n}}1 {Ñ}{{\~N}}1 {ü}{{\"u}}1 {¿}{{?`}}1 {¡}{{!`}}1
    {→}{{\textrightarrow{}}}2 {←}{{\textleftarrow{}}}2
    {·}{{$\cdot$}}1,
}
\lstset{style=wcagcode}
\usepackage[most]{tcolorbox}
\tcbuselibrary{skins, breakable}
\newtcolorbox{princbox}[2][royalblue]{enhanced, breakable,
  colback=#1!5, colframe=#1,
  fonttitle=\bfseries\normalsize, title={#2},
  arc=6pt, boxrule=1pt,
  left=10pt, right=10pt, top=8pt, bottom=8pt}
\newtcolorbox{evidbox}[1]{enhanced, breakable,
  colback=successgreen!5, colframe=successgreen,
  fonttitle=\bfseries\small, title={Evidencia: #1},
  arc=4pt, boxrule=0.8pt,
  left=8pt, right=8pt, top=6pt, bottom=6pt}
\newtcolorbox{okbox}[1][]{enhanced, breakable,
  colback=successgreen!8, colframe=successgreen!70!black,
  fonttitle=\bfseries\small, title={#1}, arc=4pt, boxrule=0.8pt,
  left=8pt, right=8pt, top=6pt, bottom=6pt}
\newtcolorbox{warnbox}[1][]{enhanced, breakable,
  colback=warningyellow!8, colframe=warningyellow,
  fonttitle=\bfseries\small, title={#1},
  arc=4pt, boxrule=0.8pt,
  left=8pt, right=8pt, top=6pt, bottom=6pt}
\usepackage[colorlinks=true, linkcolor=royalblue, urlcolor=royalblue, unicode=true,
  pdftitle={VitaPrevent - Principios WCAG y Evidencias POUR},
  pdfauthor={Erick Costa, Jonathan Tipan, Javier Quilumba},
  pdflang={es-EC}]{hyperref}
\usepackage{amssymb}
\newcommand{\cmark}{\textcolor{successgreen}{$\checkmark$}}
\newcommand{\wmark}{\textcolor{warningyellow}{$\triangle$}}
\usepackage{fancyhdr}
\pagestyle{fancy}
\fancyhf{}
\fancyhead[L]{\small\color{graymid}\textit{VitaPrevent --- Principios WCAG 2.2 y Evidencias POUR}}
\fancyhead[R]{\small\color{graymid}\thepage}
\fancyfoot[C]{\small\color{graymid}\textit{Escuela Politécnica Nacional --- 2026}}
\renewcommand{\headrulewidth}{0.4pt}
\usepackage{titlesec}
\titleformat{\section}{\normalfont\large\bfseries\color{navydeep}}{\thesection}{0.6em}{}[{\color{royalblue}\titlerule[1pt]}]
\titleformat{\subsection}{\normalfont\normalsize\bfseries\color{royalblue}}{\thesubsection}{0.6em}{}
\usepackage{parskip}
\setlength{\parindent}{0pt}
\setlength{\parskip}{6pt}
\usepackage{float}
\usepackage{setspace}
\setstretch{1.2}

% ============================================================
\begin{document}

\begin{titlepage}
\centering
\vspace*{2cm}
{\Large\bfseries\color{navydeep} ESCUELA POLITÉCNICA NACIONAL}\\[4pt]
{\normalsize Facultad de Ingeniería en Sistemas --- Ingeniería de Software}\\[4pt]
{\normalsize Asignatura: Usabilidad y Accesibilidad}\\[2cm]
{\color{royalblue}\rule{0.8\linewidth}{1.5pt}}\\[1cm]
{\Huge\bfseries\color{navydeep} VitaPrevent}\\[0.5cm]
{\Large\color{royalblue} Principios WCAG 2.2 y Evidencias de Cumplimiento}\\[0.3cm]
{\large\itshape Análisis POUR con fragmentos de código y resultados de prueba}\\[1cm]
{\color{royalblue}\rule{0.8\linewidth}{0.8pt}}\\[1.5cm]
\begin{tabular}{ll}
  \textbf{Autores:} & Erick Costa, Jonathan Tipán, Javier Quilumba \\
  \textbf{Repositorio:} & \texttt{github.com/zeuspyEC/VitaPrevent} \\
  \textbf{Sitio web:} & \texttt{vitaprevent-b2e34.web.app} \\
  \textbf{Estándar evaluado:} & WCAG 2.2 Nivel AA (W3C, octubre 2023) \\
  \textbf{Fecha:} & Junio 2026 \\
\end{tabular}
\end{titlepage}

\tableofcontents
\newpage

% ============================================================
\section{Introducción: Los principios POUR}

WCAG 2.2 organiza todos sus criterios de éxito bajo cuatro principios fundamentales, conocidos por el acrónimo \textbf{POUR}:

\begin{center}
\begin{tabular}{@{}>{\bfseries\color{navydeep}}p{2.8cm} p{11cm}@{}}
\textbf{P} --- Perceptible & El contenido debe presentarse de formas que todos los usuarios puedan percibir, incluyendo quienes no pueden ver, oír o distinguir colores. \\[6pt]
\textbf{O} --- Operable & Los componentes de interfaz y la navegación deben ser operables por cualquier usuario, independientemente de si usa ratón, teclado, voz u otras tecnologías. \\[6pt]
\textbf{U} --- Comprensible & La información y la operación de la interfaz deben ser comprensibles: el lenguaje debe ser claro, los errores descriptivos y el comportamiento predecible. \\[6pt]
\textbf{R} --- Robusto & El contenido debe ser suficientemente robusto para ser interpretado por tecnologías asistivas actuales y futuras, usando marcado válido y semántico. \\
\end{tabular}
\end{center}

VitaPrevent adopta el enfoque \textit{accessibility-first}: cada componente se diseña desde el principio para cumplir POUR, sin añadir accesibilidad como capa posterior.

% ============================================================
\section{P --- Perceptible}

\begin{princbox}[royalblue]{Principio 1: Perceptible --- La información debe poder percibirse por distintos medios}

El contenido no puede depender de un único sentido (visión, audición) para ser comprendido.
\end{princbox}

\subsection{1.1.1 --- Contenido no textual (Nivel A)}

\textbf{Criterio:} Todo contenido no textual que se presenta al usuario tiene una alternativa textual que sirve el mismo propósito.

\textbf{Cómo se cumple en VitaPrevent:}

\begin{itemize}
  \item \textbf{Emojis decorativos:} Todos los emojis usados como íconos visuales tienen \texttt{aria-hidden="true"} para que los lectores de pantalla los ignoren.
  \item \textbf{Imágenes de artículos:} El campo \texttt{imagen.alt} es obligatorio en Firestore. El componente \texttt{ArticleCard} lo consume directamente.
  \item \textbf{Cubo 3D:} Marcado con \texttt{aria-hidden="true"} --- el contenido es accesible vía Navbar.
\end{itemize}

\begin{evidbox}[1.1.1 --- Código ArticleCard con alt obligatorio]
\begin{lstlisting}[language=Java]
// ArticleCard.jsx
<img
  src={articulo.imagen?.url}
  alt={articulo.imagen?.alt || articulo.titulo}
  className="article-card__img"
  loading="lazy"
/>
\end{lstlisting}
\begin{lstlisting}[language=Java]
// Emoji decorativo en SectionHero
<span className="section-hero__icon" aria-hidden="true">
  {icon}
</span>
\end{lstlisting}
\end{evidbox}

\subsection{1.3.1 --- Información y relaciones (Nivel A)}

\textbf{Criterio:} La información, estructura y relaciones percibidas visualmente pueden determinarse programáticamente o están disponibles en texto.

\textbf{Cómo se cumple:}

\begin{evidbox}[1.3.1 --- Estructura semántica de cada página]
\begin{lstlisting}[language=XML]
<body>
  <a href="#main-content" class="skip-link">
    Saltar al contenido principal
  </a>
  <header role="banner">
    <nav aria-label="Navegación principal">...</nav>
  </header>
  <main id="main-content" tabindex="-1"
        aria-label="Contenido principal">
    <nav aria-label="Ruta de navegación"><!-- Breadcrumb --></nav>
    <section aria-labelledby="titulo-seccion">
      <h2 id="titulo-seccion">...</h2>
    </section>
  </main>
  <footer role="contentinfo">...</footer>
</body>
\end{lstlisting}
\end{evidbox}

\subsection{1.4.3 --- Contraste mínimo (Nivel AA)}

\textbf{Criterio:} La presentación visual del texto tiene una relación de contraste de al menos 4.5:1 para texto normal y 3:1 para texto grande.

\textbf{Ratios de contraste verificados con Colour Contrast Analyser:}

\begin{table}[H]
\centering
\caption{Verificación de contraste en pares de colores principales}
\begin{tabular}{@{}>{\bfseries}p{3cm} p{3cm} p{3cm} >{\centering\arraybackslash}p{2cm} >{\centering\arraybackslash}p{2cm}@{}}
\toprule
\rowcolor{royalblue!12}
\textbf{Elemento} & \textbf{Color texto} & \textbf{Color fondo} & \textbf{Ratio} & \textbf{¿Cumple?} \\
\midrule
Texto principal & \#FFFFFF blanco & \#0D1B2A navy-900 & 17.9:1 & \cmark~AA \\
Texto body & \#1F2937 gray-800 & \#FFFFFF blanco & 12.6:1 & \cmark~AA \\
Botón primario & \#FFFFFF blanco & \#1565C0 royalblue & 7.1:1 & \cmark~AA \\
Badge éxito & \#FFFFFF blanco & \#16A34A green & 4.6:1 & \cmark~AA \\
Texto hint & \#6B7280 gray-500 & \#FFFFFF blanco & 4.6:1 & \cmark~AA \\
Link en navbar & \#90CAF9 lightblue & \#0D1B2A navy-900 & 8.2:1 & \cmark~AA \\
\bottomrule
\end{tabular}
\end{table}

\subsection{1.4.10 --- Reajuste (Nivel AA)}

\textbf{Criterio:} El contenido no pierde información o funcionalidad con una anchura de 320px CSS sin scroll horizontal.

\begin{evidbox}[1.4.10 --- Sistema responsive con tokens CSS]
\begin{lstlisting}[language=bash]
/* tokens.css */
--breakpoint-sm: 480px;
--breakpoint-md: 768px;
--breakpoint-lg: 1024px;

/* Home.css */
@media (min-width: 1024px) {
  .hero__content {
    grid-template-columns: 1fr 1fr;
  }
}
/* En mobile (< 480px): una columna, sin scroll horizontal */
\end{lstlisting}
\end{evidbox}

% ============================================================
\section{O --- Operable}

\begin{princbox}[successgreen]{Principio 2: Operable --- La interfaz debe poder utilizarse con distintos mecanismos}

Ninguna funcionalidad debe requerir exclusivamente el uso del ratón o gestos complejos.
\end{princbox}

\subsection{2.1.1 --- Teclado (Nivel A)}

\textbf{Criterio:} Toda la funcionalidad del contenido es operable mediante una interfaz de teclado, sin requerir tiempos específicos para cada pulsación.

\textbf{Verificación de ruta Tab en VitaPrevent:}

\begin{center}
\begin{tabular}{>{\bfseries\small}r l}
1 & SkipLink (``Saltar al contenido principal'') \\
2 & Logo VitaPrevent (enlace a Inicio) \\
3--9 & Ítems de Navbar (Inicio, Atención Primaria, Vacunación, \ldots) \\
10 & Primer enlace/botón interactivo en el contenido principal \\
$\vdots$ & \ldots (continúa en orden lógico visual) \\
\end{tabular}
\end{center}

Todos los filtros, botones de acordeón, campos de formulario y botones de paginación son alcanzables y accionables solo con teclado.

\subsection{2.1.2 --- Sin trampa de teclado (Nivel A)}

\textbf{Criterio:} Si el foco del teclado se puede mover a un componente de la página mediante una interfaz de teclado, el foco puede también alejarse de ese componente usando solo teclado.

\begin{evidbox}[2.1.2 --- Modal con focus trap y Escape]
\begin{lstlisting}[language=Java]
// hooks/useFocusTrap.js
useEffect(() => {
  const focusables = modal.current
    .querySelectorAll('button, [href], input, select, [tabindex]')
  const first = focusables[0]
  const last  = focusables[focusables.length - 1]

  const trap = (e) => {
    if (e.key !== 'Tab') return
    if (e.shiftKey && document.activeElement === first) {
      e.preventDefault(); last.focus()
    } else if (!e.shiftKey && document.activeElement === last) {
      e.preventDefault(); first.focus()
    }
  }
  modal.current.addEventListener('keydown', trap)
  return () => modal.current?.removeEventListener('keydown', trap)
}, [])
\end{lstlisting}
\end{evidbox}

\subsection{2.4.1 --- Evitar bloques (Nivel A)}

\textbf{Criterio:} Existe un mecanismo para evitar los bloques de contenido que se repiten en múltiples páginas.

\begin{evidbox}[2.4.1 --- SkipLink como primer elemento del DOM]
\begin{lstlisting}[language=Java]
// SkipLink.jsx - siempre el PRIMER elemento del DOM
export default function SkipLink() {
  return (
    <a href="#main-content" className="skip-link">
      Saltar al contenido principal
    </a>
  )
}
// SkipLink.css
.skip-link {
  position: fixed;
  transform: translateY(-200%);  /* oculto por defecto */
}
.skip-link:focus-visible {
  transform: translateY(0);      /* visible al recibir Tab */
  outline: 3px solid #1e88e5;
}
\end{lstlisting}
\end{evidbox}

\subsection{2.4.3 --- Orden de foco (Nivel A)}

\textbf{Criterio:} Si una página web puede navegarse secuencialmente, los componentes que reciben el foco lo hacen en un orden que preserva el significado y la operabilidad.

\textbf{Bug encontrado y corregido:} \texttt{PageWrapper} movía el foco al \texttt{<main>} en cada carga de página, incluyendo la carga inicial. Esto provocaba que Tab comenzara en el primer botón del contenido, saltando el SkipLink y la Navbar.

\begin{evidbox}[2.4.3 --- Corrección en PageWrapper.jsx]
\begin{lstlisting}[language=Java]
// ANTES (incorrecto): focus() en carga inicial
useEffect(() => {
  document.title = pageTitle
  mainRef.current?.focus()  // saltaba Navbar en carga inicial
}, [pathname, title])

// DESPUÉS (correcto): solo en navegación SPA
const isFirstRender = useRef(true)
useEffect(() => {
  document.title = pageTitle
  if (isFirstRender.current) {
    isFirstRender.current = false
    return   // carga inicial: Tab va SkipLink -> Navbar
  }
  mainRef.current?.focus()  // SPA: anuncia nuevo contenido
}, [pathname, title])
\end{lstlisting}
\end{evidbox}

\subsection{2.4.7 --- Foco visible (Nivel AA)}

\textbf{Criterio:} Cualquier interfaz de teclado tiene un modo de operación en el que el indicador de foco del teclado es visible.

\begin{evidbox}[2.4.7 --- Focus visible con CSS personalizado]
\begin{lstlisting}[language=bash]
/* global.css */
:focus-visible {
  outline: 3px solid var(--color-focus, #1e88e5);
  outline-offset: 3px;
  border-radius: 4px;
}
/* Elimina el outline en click (solo en navegación por teclado) */
:focus:not(:focus-visible) {
  outline: none;
}
\end{lstlisting}
\end{evidbox}

\subsection{2.4.11 --- Apariencia del foco (Nivel AA, WCAG 2.2)}

\textbf{Criterio (nuevo en WCAG 2.2):} El indicador de foco del teclado es de al menos el área del perímetro del componente, tiene relación de contraste de al menos 3:1 entre los píxeles enfocados y no enfocados, y no está completamente oculto por otro contenido.

\textbf{Verificación:} El outline de 3px \#1e88e5 sobre fondo blanco tiene ratio 3.2:1. Se aplica sobre el borde exterior del componente, no oculto por nada.

% ============================================================
\section{U --- Comprensible}

\begin{princbox}[purplemod]{Principio 3: Comprensible --- El contenido debe ser claro y predecible}

Los usuarios deben poder entender tanto el contenido como el funcionamiento de la interfaz.
\end{princbox}

\subsection{3.1.1 --- Idioma de la página (Nivel A)}

\textbf{Criterio:} El idioma humano predeterminado de cada página web puede determinarse programáticamente.

\begin{evidbox}[3.1.1 --- Lang en index.html]
\begin{lstlisting}[language=XML]
<!-- index.html -->
<!DOCTYPE html>
<html lang="es">
  <head>
    <meta charset="UTF-8" />
    <title>VitaPrevent</title>
  </head>
  ...
</html>
\end{lstlisting}
\end{evidbox}

\subsection{3.3.1 y 3.3.2 --- Identificación y etiquetas de error (Nivel A)}

\textbf{Criterio 3.3.1:} Si se detecta automáticamente un error en la entrada, el elemento en error se identifica y el error se describe al usuario en texto.

\textbf{Criterio 3.3.2:} Si el contenido requiere entrada del usuario, se proporcionan etiquetas o instrucciones.

\begin{evidbox}[3.3.1 / 3.3.2 --- Componente FormField accesible]
\begin{lstlisting}[language=Java]
// FormField.jsx
<div className={`form-field ${error ? 'form-field--error' : ''}`}>
  <label htmlFor={id} className="form-field__label">
    {label}
    {required && <span aria-hidden="true"> *</span>}
    {required && <span className="sr-only">(requerido)</span>}
  </label>

  {hint && <span id={`${id}-hint`} className="form-field__hint">
    {hint}
  </span>}

  <input
    id={id}
    aria-required={required}
    aria-invalid={!!error}
    aria-describedby={
      [hint && `${id}-hint`, error && `${id}-error`]
        .filter(Boolean).join(' ')
    }
  />

  {error && (
    <span id={`${id}-error`} className="form-field__error"
          role="alert">
      {error}
    </span>
  )}
</div>
\end{lstlisting}
\end{evidbox}

\subsection{3.2.3 / 3.2.4 --- Navegación y identificación coherentes (Nivel AA)}

\textbf{Criterio:} Los mecanismos de navegación que se repiten en múltiples páginas aparecen en el mismo orden relativo. Los componentes con el mismo propósito se identifican de manera coherente.

\textbf{Verificación:} La Navbar es idéntica en todas las páginas --- mismo componente React reutilizado. \texttt{aria-current="page"} se actualiza dinámicamente según la ruta activa:

\begin{evidbox}[3.2.3 / 3.2.4 --- aria-current dinámico en Navbar]
\begin{lstlisting}[language=Java]
// Navbar.jsx
{NAV_ITEMS.map((item) => (
  <NavLink
    key={item.path}
    to={item.path}
    aria-current={isActive(item.path) ? 'page' : undefined}
    className={({ isActive }) =>
      `nav-link ${isActive ? 'nav-link--active' : ''}`
    }
  >
    {item.label}
  </NavLink>
))}
\end{lstlisting}
\end{evidbox}

% ============================================================
\section{R --- Robusto}

\begin{princbox}[errored]{Principio 4: Robusto --- El contenido debe ser compatible con tecnologías actuales y futuras}

El contenido debe poder ser interpretado de manera fiable por una amplia variedad de agentes de usuario, incluidas las tecnologías asistivas.
\end{princbox}

\subsection{4.1.1 --- Procesamiento (Nivel A)}

\textbf{Criterio:} El contenido implementado usando lenguajes de marcado tiene elementos con etiquetas de inicio y fin completas, no contiene atributos duplicados, e IDs únicas.

\textbf{Verificación:} El código JSX de React garantiza HTML cerrado correctamente. Se verificó con el W3C Validator sin errores. Los IDs de secciones son únicos por página (\texttt{id="skills-title"} vs \texttt{id="recursos-sm-title"}).

\subsection{4.1.2 --- Nombre, función, valor (Nivel A)}

\textbf{Criterio:} Para todos los componentes de interfaz de usuario, el nombre y función pueden determinarse programáticamente; los estados, propiedades y valores que el usuario puede establecer pueden establecerse programáticamente.

\begin{evidbox}[4.1.2 --- Accordion con ARIA completo]
\begin{lstlisting}[language=Java]
// Accordion.jsx
<div className="accordion-item" key={item.id}>
  <h3 className="accordion-item__heading">
    <button
      id={`accordion-trigger-${item.id}`}
      aria-expanded={openIds.has(item.id)}
      aria-controls={`accordion-panel-${item.id}`}
      onClick={() => toggle(item.id)}
      className="accordion-item__trigger"
    >
      {item.question}
      <span aria-hidden="true" className="accordion-item__icon">
        {openIds.has(item.id) ? '-' : '+'}
      </span>
    </button>
  </h3>
  <div
    id={`accordion-panel-${item.id}`}
    role="region"
    aria-labelledby={`accordion-trigger-${item.id}`}
    hidden={!openIds.has(item.id)}
    className="accordion-item__panel"
  >
    <p>{item.answer}</p>
  </div>
</div>
\end{lstlisting}
\end{evidbox}

\subsection{4.1.3 --- Mensajes de estado (Nivel AA)}

\textbf{Criterio:} En el contenido implementado usando lenguajes de marcado, los mensajes de estado pueden determinarse programáticamente mediante rol o propiedades.

\begin{evidbox}[4.1.3 --- aria-live en calculadora IMC y Spinner]
\begin{lstlisting}[language=Java]
// Resultado de la calculadora IMC
<div
  className={`imc-result imc-result--${imcResult.color}`}
  role="status"
  aria-live="polite"
>
  <span aria-label={`IMC: ${imcResult.valor}`}>
    {imcResult.valor}
  </span>
  <Badge variant={imcResult.color}>{imcResult.categoria}</Badge>
  <p>{imcResult.descripcion}</p>
</div>

// Spinner de carga
<div role="status" aria-label={label || 'Cargando...'}>
  <div className="spinner" aria-hidden="true" />
</div>
\end{lstlisting}
\end{evidbox}

% ============================================================
\section{Resumen: Cumplimiento POUR por principio}

\begin{table}[H]
\centering
\caption{Cumplimiento POUR global de VitaPrevent}
\begin{tabular}{@{}>{\bfseries\color{navydeep}}p{3cm} >{\centering\arraybackslash}p{2.5cm} p{8cm}@{}}
\toprule
\rowcolor{royalblue!12}
\textbf{Principio} & \textbf{Cumplimiento} & \textbf{Criterios destacados} \\
\midrule
Perceptible & \textcolor{successgreen}{\bfseries 97\%} & 1.1.1 alt text, 1.3.1 semántica, 1.4.3 contraste 17.9:1, 1.4.10 responsive 320px \\
Operable & \textcolor{successgreen}{\bfseries 93\%} & 2.1.1 teclado, 2.4.1 SkipLink, 2.4.3 Tab fix, 2.4.7 focus visible. Parcial: 2.5.7 drag \\
Comprensible & \textcolor{successgreen}{\bfseries 95\%} & 3.1.1 lang, 3.2.3 Navbar coherente, 3.3.1--2 errores descriptivos. Parcial: 3.2.6 ayuda \\
Robusto & \textcolor{successgreen}{\bfseries 100\%} & 4.1.1 HTML válido, 4.1.2 ARIA completo, 4.1.3 aria-live \\
\midrule
\textbf{Global WCAG 2.2 AA} & \textcolor{successgreen}{\bfseries 96\%} & 2 criterios con cumplimiento parcial documentado \\
\bottomrule
\end{tabular}
\end{table}

\begin{okbox}[Declaración resumida de conformidad]
VitaPrevent cumple sustancialmente con los cuatro principios POUR de WCAG 2.2 Nivel AA. Los dos criterios con cumplimiento parcial (2.5.7 y 3.2.6) tienen alternativas funcionales documentadas y están en proceso de resolución completa. El sitio \textbf{no presenta barreras de acceso} para usuarios con discapacidades visuales, motrices, auditivas o cognitivas que utilizan lectores de pantalla, teclado o magnificadores.
\end{okbox}

\end{document}
